\documentclass[12pt,a4paper]{article}

\usepackage{cmap}
\usepackage[T2A]{fontenc}
\usepackage[utf8]{inputenc}
\usepackage[russian]{babel}
\usepackage{amsmath,amssymb}
\usepackage[a4paper,margin=2cm]{geometry}

\setlength{\parindent}{0pt}
\setlength{\parskip}{0.45em}
\sloppy

\begin{document}

\begin{center}
{\Large\bfseries Билет 3: вопросы для проверки знаний}
\end{center}

\noindent\fbox{%
  \begin{minipage}{\dimexpr\linewidth-2\fboxsep-2\fboxrule\relax}
  \normalfont
Частные производные функции нескольких переменных. Дифференцируемость
функции в точке, дифференциал. Необходимые условия дифференцируемости,
достаточные условия дифференцируемости функции нескольких переменных.
Дифференцируемость сложной функции. Инвариантность формы дифференциала
относительно замены переменных. Градиент и производная по направлению.
  \end{minipage}%
}

\section*{Вопросы на понимание}

\begin{enumerate}
\item Чем частная производная по \(x_i\) отличается от производной по произвольному вектору \(\ell\), если обе вычисляются через пределы функций одной переменной?
\item Почему при сравнении скоростей роста функции рассматривают производные по единичным направлениям, а не по произвольным ненормированным векторам?
\item Что означает, что \(df(x^0)\) является главной линейной частью приращения \(\Delta f\)? Когда \(\Delta f\) нельзя просто заменить на \(df\) без учета остатка?
\item Почему условие остатка \(o(|x-x^0|)\) в определении дифференцируемости сильнее, чем проверка приращений только вдоль координатных осей?
\item Для функции
\[
  f(x,y)=
  \begin{cases}
  \dfrac{xy}{x^2+y^2}, & (x,y)\ne(0,0),\\
  0, & (x,y)=(0,0),
  \end{cases}
\]
объясните, почему существование частных производных в \((0,0)\) еще не означает непрерывность и дифференцируемость.
\item Почему дифференцируемость функции в точке автоматически дает непрерывность в этой точке? Почему непрерывность сама по себе не должна давать дифференцируемость?
\item Зачем в достаточном условии дифференцируемости требуется, чтобы частные производные были определены в окрестности точки и непрерывны в самой точке?
\item Чем достаточное условие дифференцируемости отличается от необходимого? Можно ли по нему описать все дифференцируемые функции?
\item Если функция дифференцируема, какую информацию о производных по направлениям содержит \(\nabla f(x^0)\)? Что меняется, если \(\nabla f(x^0)=0\)?
\item Почему направление градиента является направлением наибольшего роста только при \(\nabla f(x^0)\ne0\)?
\item В правиле цепочки \(\operatorname{Jac}(G\circ F)(x^0)=\operatorname{Jac}G(F(x^0))\operatorname{Jac}F(x^0)\) как размерности матриц объясняют порядок умножения?
\item В чем смысл инвариантности формы первого дифференциала при замене переменных: что остается неизменным в записи, а что меняется в смысле \(dy\)?
\end{enumerate}

\section*{Источники для сверки}

\texttt{target/answer\_question\_03.tex};
\texttt{IvGE\_1.pdf}, стр. 180--191.

\end{document}
